\documentclass{article}
\usepackage[utf8]{inputenc}
\usepackage{amsmath}
\usepackage{geometry}
\geometry{margin=2cm}
\begin{document}

\section*{Questão 144 — ENEM 2024 — Matemática}

\subsection*{Enunciado}
Para obter um sólido de revolução (rotação de 360$^\circ$ em torno de um eixo fixo), uma professora realizou as seguintes etapas:
\begin{itemize}
  \item Recortou o trapézio retângulo PQRS de um material rígido;
  \item Afixou o lado PS do trapézio em uma vareta fixa retilínea (eixo de rotação);
  \item Girou o trapézio 360$^\circ$ em torno da vareta e obteve um sólido de revolução.
\end{itemize}
Observe a figura que apresenta o trapézio afixado na vareta e o sentido de giro.

\vspace{0.5em}
\textbf{O sólido obtido foi um(a):}

\subsection*{Solução Técnica}
A alternativa correta é a letra \textbf{D}, pois a figura gerada pela rotação do trapézio retângulo em torno do lado fixado forma um \textbf{tronco de cone}.

\subsection*{Estratégia e Passos}
Identificou-se o sólido gerado pela rotação do trapézio PQRS em torno do lado PS fixado:
\begin{enumerate}
  \item Trapézio retângulo com ângulos retos em P e S; lado PS perpendicular às bases PQ e RS.
  \item PS fixo como eixo de rotação.
  \item Ao girar o trapézio 360$^\circ$ em torno de PS, cada ponto desloca-se em círculo perpendicular ao eixo.
  \item As bases PQ e RS geram duas circunferências de raios diferentes.
  \item O sólido resultante tem duas bases circulares paralelas e raios diferentes, ligadas por superfície lateral inclinada.
  \item Este sólido é característico de um tronco de cone.
\end{enumerate}

\subsection*{Explicação Didática}
Ao girar o trapézio em torno do lado PS, criamos uma figura tridimensional com as características de um tronco de cone. Isso ocorre porque as bases do trapézio geram circunferências de diferentes raios, resultando em bases paralelas circulares diferentes, não iguais como no cilindro, e tampouco pontiagudas como no cone.

\subsection*{Erros Comuns e Distratores}
Foram identificadas opções distractoras baseadas em confusões comuns:
\begin{itemize}
  \item \textbf{A}: Pensar que o sólido é um cone, por confundir trapézio retângulo com triângulo.
  \item \textbf{B}: Supor que o sólido é um cilindro, ignorando os comprimentos diferentes das bases.
  \item \textbf{C}: Associar o sólido a uma pirâmide, metonímia entre sólido de revolução e sólido poligonal.
  \item \textbf{E}: Confundir tronco de pirâmide com tronco de cone, misturando bases poligonais com circulares.
\end{itemize}

\subsection*{Perfil da Questão}
\begin{itemize}
  \item Habilidades principais: Geometria espacial e interpretação de sólidos de revolução.
  \item Demanda cognitiva média, com requisito de visualização espacial e noção de figuras tridimensionais.
  \item Não exige cálculos, mas compreensão conceitual da formação do sólido pela revolução.
\end{itemize}

\subsection*{Conclusão}
O sólido obtido pela rotação do trapézio retângulo em torno do lado PS fixo é um \textbf{tronco de cone}. Esta conclusão é suportada pela observação dos deslocamentos das bases do trapézio e pela análise da superfície lateral resultante da rotação.

\vspace{1em}
\noindent\textit{Figura:} Trapézio retângulo PQRS fixado no lado PS que gira em torno da vareta, gerando o sólido de revolução.

\vspace{1em}
\noindent\textbf{Resposta correta: tronco de cone (letra D).

\end{document}